A genuine perception released lightly into the interaction and not 
pursued beyond the point of contact. The thought is sufficiently 
formed to stand on its own, but it is not extended, explained, or 
resolved --- and crucially, it is not lingered in. The speaker 
moves on. The meaning stays with her.

This is what distinguishes Graze from a placement that simply 
stops. A graze passes through the moment. The linger lives in the 
receiver, not in the one who expressed it. That is what preserves 
the space --- she is left holding something he has already moved 
past, which means her completion of it is entirely her own.

Because a graze releases something real at low intensity and then 
moves on, it functions as the clearest diagnostic of the three 
modes --- what she does next is uncontaminated by pressure or 
momentum. This makes Graze the most calibrated expression in the 
ladder, and the natural entry point into co-authorship.

It is not a technique for creating openness. It is what genuine 
perception looks like at its lightest register --- touched, not 
pursued. Passed through, not held.

